\documentclass[12pt]{article}
\usepackage[utf8]{inputenc}
\usepackage[T1]{fontenc}
\usepackage[margin=1in]{geometry}
\usepackage{setspace}
\usepackage{xcolor}
\usepackage[colorlinks=true, urlcolor=blue!60!black, linkcolor=blue!60!black]{hyperref}
\setlength{\parindent}{0em}
\setlength{\parskip}{0.6em}
\renewcommand{\baselinestretch}{1.0}
\begin{document}
\begin{center}
{\LARGE \textit{Egbert v. Boule}, 596 U.S. 482 (2022)\par}
\vspace{0.5em}
{\large Case Brief --- Two-Page Edition\par}
\end{center}

\section{Memory Jogger}
An innkeeper roughed up by a Border Patrol agent could not sue under \textit{Bivens}: no remedy if there is any reason to think Congress is better equipped to create one (\href{https://supreme.justia.com/cases/federal/us/596/21-147/}{Justia}).

\section{Detailed Case Facts}
Robert Boule's Smuggler's Inn abutted the Canadian border. Border Patrol Agent Egbert, following a guest into the inn's driveway, allegedly threw Boule to the ground when asked to leave, then retaliated against Boule's complaints by prompting an IRS audit.

\section{Procedural History}
The district court granted summary judgment for Egbert; the Ninth Circuit reversed, allowing both \textit{Bivens} claims.

\section{Judicial Votes}
Reversed: 6--3 on the Fourth Amendment claim, 9--0 on the First. Thomas wrote for Roberts, Alito, Kavanaugh, and Barrett; Gorsuch concurred in the judgment; Sotomayor (with Breyer and Kagan) concurred as to the First Amendment and dissented as to the Fourth.

\section{Holding}
\textit{Bivens} provides no damages remedy for either claim. A court may not recognize a \textit{Bivens} action if there is \textit{any} reason to think Congress might be better equipped to create a damages remedy.

\section{Analysis of Opinions}
\textbf{Majority:} Creating damages actions is legislative work. Border security and CBP's internal grievance process each foreclosed the excessive-force claim, and retaliation claims would chill legitimate conduct across the federal workforce. \textbf{Gorsuch:} the Court should overrule \textit{Bivens} outright. \textbf{Sotomayor:} the force claim was ``materially indistinguishable'' from \textit{Bivens}, and the grievance process offered no remedy at all.

\section{Examples: Future Applications}
\textbf{Same-Side:} a First Amendment retaliation suit against a federal inspector fails. \textbf{Opposite-Side:} the same force by a state trooper is actionable under 42 U.S.C. \S~1983, which Congress created. \textbf{Fence-Sitter:} federal narcotics agents on a joint task force conduct a warrantless arrest and strip search that closely tracks \textit{Bivens} itself.

\section{Critique}
State officers answer in damages under \S~1983 for conduct that federal officers now largely escape, and any internal complaint process can now immunize federal agents. Justice Gorsuch faulted the Court for keeping \textit{Bivens} in name while making it nearly impossible to invoke.

\section{Key Quotations}
\begin{quote}
``While our cases describe two steps, those steps often resolve to a single question: whether there is any reason to think that Congress might be better equipped to create a damages remedy.''
\end{quote}
\noindent --- Justice Thomas, for the Court, 596 U.S. at 492.

\end{document}
